\section{所有权}
\begin{enumerate}
  \item 必考。
  \item \zd{区分原则，物债两分}。
  \item 公示方法。
        \begin{enumerate}
          \item 不动产登记。\zd{登记生效主义}。
                \begin{enumerate}
                  \item 更正登记 vs 异议登记。请求更正登记 \yjt 权利人不同意更正，权属有争议。\al{借名买房}。
                        \begin{enumerate}
                          \item 15 日内去起诉。
                          \item 不阻却债权、也不阻却物权效力。
                          \item 第三人无法善意取得。
                        \end{enumerate}
                  \item \zd{预告登记}。较重要。从债权到物权有风险。\zd{不阻却债权，但未经同意不发生所有权变动}。
                        \begin{enumerate}
                          \item 能够登记日 90 日之内。\al{一房数卖；期房、二手房}。
                        \end{enumerate}
                \end{enumerate}
          \item 动产交付。
                \begin{enumerate}
                  \item 现实交付。主观上有移转标的物所有权的意思，客观上完成占有的变动。
                  \item \zd{观念交付}。必考，一个选项。为了减少交易成本、风险，提高效率。
                        \begin{enumerate}
                          \item 简易交付。先占有（有权 \hz 无权）后购买，买卖合同\zd{生效时}。
                          \item 指示交付。\al{甲出租车给乙，后出国卖给丙}。
                                \begin{itemize}
                                  \item 3 方关系，\zd{必然发生了标的物的所有人和占有人分离}。
                                  \item 典型：自己所有他人占有。本质是\zd{转让返还原物请求权}（返还请求权的让与）。
                                  \item 买卖合同\zd{生效时}，内部 \yjt 甲丙之间；外部 \yjt 通知乙。
                                \end{itemize}
                          \item 占有改定。\zd{所有人变占有人}，先卖后租/借；\zd{租赁/借用合同生效时}。融资租赁合同。
                        \end{enumerate}
                \end{enumerate}
        \end{enumerate}
  \item 物权变动。
        \begin{enumerate}
          \item 一般的物权变动。
                \begin{enumerate}
                  \item 基于民事法律行为 \hz 合同（买卖 \hz 赠予）。\zd{有权处分 + 有效合同 + 登记/交付 = 继受取得}。
                  \item 非基于民事法律行为。不需要登记、交付。
                        \begin{enumerate}
                          \item 政府征收，征收决定生效时。
                          \item 法律文书，常考点，法院 \hz 仲裁文书生效时（\zd{形成}）。
                          \item 唯一继承，被继承人死亡时。
                          \item 事实行为，合法建造成就时/完成时。
                        \end{enumerate}
                \end{enumerate}
          \item 特殊\blh
                \begin{enumerate}
                  \item \zd{善意取得 = 无权处分 + 有效合同 + 登记/交付 = 原始取得}。3 方关系。
                        \begin{enumerate}
                          \item \zd{无权处分}，\zd{以自己名义}处分他人或共有财产。和有权处分可以相互转化。
                                \al{甲表乙保管，乙出卖给丙，甲请求乙还钱}。要钱不要物。对处分行为进行了\zd{追认}。
                          \item 相对人主客观相结合，善意（到完成公示的时候都不知道），无重大过失；支付合理对价。
                          \item \zd{不准用占有改定}，公示效力极弱。\al{甲委托乙保管字画，乙出卖给丙时约定先借来把玩几天}。
                          \item 救济：违约责任、侵权责任、不当得利，择一。
                          \item 动产（合法）占有委托物可以；
                          \item \zd{盗赃、埋藏、隐藏、遗失不适用善意取得，但可继受取得}，追及效力。
                                \begin{itemize}
                                  \item \al{甲表 \yjt 乙捡 \yjt\, 3 年}。催促行权，可变可中断。
                                  \item \zd{遗失物再转让}。\al{甲表 \yjt 乙捡卖给 \yjt 丙 \yjt\, 2 年}。剥夺实体权利，不变。
                                  \item \al{甲表 \yjt 乙捡卖给 \yjt 商事主体卖给 \yjt 丙}。\zd{具有经营资格的经营者或者通过拍卖的手段}。
                                \end{itemize}
                        \end{enumerate}
                  \item 动产遗失物的拾得。\zd{重点保护失主利益，而非交易秩序和安全}。
                        \begin{enumerate}
                          \item \zd{保管义务}，无故意或重大过失，否则不赔；\zd{返还义务}。
                          \item \zd{必要费用请求权}，准用无因管理。
                          \item \zd{悬赏报酬请求}，单方允诺之债。
                          \item \zd{不能侵占，否则丧失以上各种请求权}；1 年后归国家。
                        \end{enumerate}
                  \item 先占。动产、无主物。事实行为，原始取得。
                  \item 添附。\zd{不同人的}物结合。附和（\al{油漆柜子、木板钉子}）、混合（\al{油和油}）、加工（动产和劳动）。
                \end{enumerate}
        \end{enumerate}
  \item 物权变动的学理分类。
        \begin{enumerate}
          \item 原始取得（物的所有权的第一次产生）\yjt 法定；
                \begin{enumerate}
                  \item 从无到有：先占、孳息、劳动生产。
                  \item 从有到有：善意取得、拾得、发现、添附、征收。
                \end{enumerate}
          \item \zd{继受取得} \yjt 意定。买卖/赠与合同；继承 \hz 受遗赠。
        \end{enumerate}
  \item 业主的建筑区分所有权。
        \begin{enumerate}
          \item 对房屋的套内空间面积享有专用权。
          \item 共有部分的共有权，电梯、楼道、消防设施、\zd{承重墙}。不得以放弃权利为由不履行义务，\al{物业费}。
                \begin{enumerate}
                  \item 车位。
                        \begin{enumerate}
                          \item 原始。开发商，卖、租、赠。
                          \item 新增。全小区业主。
                        \end{enumerate}
                  \item 利用业主的共有部分产生的收入。
                  \item 墙体中间层一人一半。\al{装修钻透漏板；漏水}。
                \end{enumerate}
          \item 管理权、成员权、表决权。
        \end{enumerate}
  \item 共有。
        \begin{enumerate}
          \item 按份共有。人合性，人身信赖关系。\zd{内外有别}。
                \begin{enumerate}
                  \item 内部，有约从约。
                        \begin{enumerate}
                          \item 保存行为。保值、增值行为。
                          \item 重大事项。重大修缮 \yjt\, 2/3、变更性质或用途。
                          \item \zd{出租行为}。单独的负担行为。
                          \item 收益分配。有约从约；按出资额；份额确定（平均分）。
                          \item 费用负担。同上。
                          \item 分割。
                          \item 处分。2/3。\zd{卖东西 \yjt 物权变动（继受 \hz 善意取得）；卖份额 \yjt 优先购买权}。
                        \end{enumerate}
                  \item 优先购买权。\zd{对外转让、同等条件、转让自由}。
                        \begin{enumerate}
                          \item 对外部而言，对内意味着破坏人合性。\al{甲乙丙共有手表，乙 10\% 份额转让给丙}。
                          \item \zd{转让自由}，通知（诚信原则），15 天。\al{乙 10\% 份额转让给丁}。
                          \item \zd{同等规则/镜面规则}。都要买，有约从约；协商；按比例行权（不影响原来架构）。
                        \end{enumerate}
                  \item 外部（侵权问题）。共同侵权，对外连带，对内按份。
                \end{enumerate}
          \item 共同共有。夫妻、家庭、遗产继承、合伙关系、同居。
        \end{enumerate}
\end{enumerate}
